\bookchapter{Глава 57. Путь дальше — разрешения, уведомления и современная архитектура}

\chapterintro{Уметь собрать приложение — только начало. В этой главе поговорим о ступенях от «умею делать» до «делаю хорошо»: как изящно просить разрешения, как слать уведомления, не раздражая людей, как поднять архитектуру ещё на этаж и — как в 2026 году писать Android-код в паре с ИИ. Каждая из этих тем встанет перед вами в настоящем проекте очень скоро.}

\booksection{Разрешения во время выполнения: сначала постучались, потом вошли}

Камера, геолокация, контакты, уведомления… Для этих «чувствительных зон» Android велит спрашивать у пользователя прямо во время работы приложения (начиная с Android 6). Порядок такой: сначала объявить (Manifest), потом запросить (код), а на отказ пользователя — изящно обойтись без неё. В Compose проще всего через rememberPermissionState из официальной библиотеки Accompanist:

\begin{lstlisting}[language=Kotlin]
// Сначала объявляем в манифесте:
// <uses-permission android:name="android.permission.CAMERA" />

@Composable
fun CameraButton() {
    val permission = rememberPermissionState(
        permission = android.Manifest.permission.CAMERA
    )
    Button(onClick = {
        when {
            permission.status.isGranted -> openCamera()
            permission.status.shouldShowRationale -> showExplainDialog()
            else -> permission.launchPermissionRequest()
        }
    }) { Text("Сделать фото") }
}
\end{lstlisting}

\begin{tcbwarn}{Внимание: а если отказали навсегда?}
Если пользователь отметил «больше не спрашивать», launchPermissionRequest больше не покажет окно. Тогда мягко направьте пользователя в системные настройки, чтобы включить разрешение вручную (перейдите туда через Intent на страницу настроек приложения), и объясните на экране, «зачем приложению это разрешение». Чем вежливее тон, тем охотнее пользователь даёт разрешение — текст запроса разрешения стоит написать с душой.

\end{tcbwarn}

\booksection{Уведомления: в нужный момент — полезные слова}

Уведомления — канал, по которому приложение достаёт пользователя, и одновременно функция, которую первыми «банят». Три дисциплины уведомлений: только полезное, возможность отключить, никаких ночных налётов. Базовое применение:

\begin{lstlisting}[language=Kotlin]
// На Android 13+ сначала нужно запросить разрешение POST_NOTIFICATIONS (см. предыдущий раздел)

// 1. Создаём канал уведомлений (обязательно на Android 8+)
val channel = NotificationChannel(
    "reminder", "Напоминания", NotificationManager.IMPORTANCE_DEFAULT
)
val manager = getSystemService(NotificationManager::class.java)
manager.createNotificationChannel(channel)

// 2. Отправляем уведомление
val notification = NotificationCompat.Builder(this, "reminder")
    .setSmallIcon(R.drawable.ic_notification)
    .setContentTitle("Напоминание о записях")
    .setContentText("Сегодня ещё нет записей")
    .setPriority(NotificationCompat.PRIORITY_DEFAULT)
    .build()
if (ActivityCompat.checkSelfPermission(this,
        Manifest.permission.POST_NOTIFICATIONS) == PackageManager.PERMISSION_GRANTED
) {
    manager.notify(1, notification)
}
\end{lstlisting}

Отложенные задачи (например, «каждый день в 20:00 напоминать записать траты») делайте через WorkManager — официально рекомендованного «дворецкого фоновых задач»: он гарантирует выполнение в подходящий момент и избавляет от возни с оптимизацией батареи. Для будильничных задач это беспроигрышный выбор.

\booksection{Архитектура на вырост: внедрение зависимостей и модульность}

Когда проект подрастёт, заявятся две «сладкие беды»:

\begin{itemize}
  \item Хаос с созданием объектов: каждый экран сам делает new для базы, new для API — наведите единый порядок через Hilt (фреймворк внедрения зависимостей): «кому что надо — попросил, и подано»;
  \item Код громоздится горой: проект в одном модуле правится всё медленнее — разбейте его на несколько модулей (app / core / feature-xxx), чтобы каждую функцию вести и тестировать независимо.
\end{itemize}

\begin{tcbinfo}{Совет: не переусложняйте}
Hilt и модульность — хорошие инструменты, но в одиночном проекте не торопитесь. Учётной книге с избытком хватает одного модуля и ручного создания зависимостей. Подключайте их по необходимости, когда код перевалит за десятки тысяч строк, компиляция начнёт тормозить и к делу присоединятся ещё люди. Инструменты служат цели, а не показухе, — это и есть наш принцип «простота превыше всего» из книги о программировании с ИИ.

\end{tcbinfo}

\booksection{Новая реальность 2026 года: Android-разработка с ИИ-напарником}

В 2026 году ИИ глубоко врос в процесс Android-разработки. Если хотите ускориться, вот проверенный набор приёмов:

\begin{center}
\begin{tabularx}{\linewidth}{|X|X|X|}
\hline
\textbf{Этап} & \textbf{Чем помогает ИИ} & \textbf{Ваша зона ответственности} \tabularnewline
\hline
Написание кода & Дополнение и генерация встроенным в Android Studio Gemini или сторонними ИИ-плагинами & Ясно описать требования и проверить результат \tabularnewline
Поиск информации & Отдать официальную документацию и текст ошибки ИИ на растолкование & Ключевые выводы сверить с официальной документацией \tabularnewline
Написание тестов & Попросить ИИ сгенерировать кейсы трёх видов: «нормальные / граничные / ошибочные» & Убедиться, что кейсы отвечают смыслу задачи \tabularnewline
Правка багов & Вставить стек и относящийся код — ИИ найдёт и починит & Разобраться в правках и проверить их тестами \tabularnewline
\hline
\end{tabularx}
\end{center}

\begin{center}\small\color{gray}Четыре сценария Android-разработки с ИИ-напарником\end{center}

Конкретная методика — как ставить ИИ задачи, как принимать работу, как обходить ловушки — в полном объёме изложена в 11 главах седьмого раздела этой книги «Практика программирования с ИИ» (главы 37–47); очень советую читать в связке: синтаксис Android — здешняя часть, а искусство сотрудничества с ИИ — седьмой раздел.

\booksection{Оптимизация производительности Compose и инструменты отладки}

Писать на Compose — удовольствие, но «слишком много рекомпозиций» тоже подтормаживает. Три практических приёма:

\begin{itemize}
  \item Стабильные параметры: добавляйте к классам данных @Immutable или используйте derivedStateOf, чтобы изменения несвязанного состояния не вызывали масштабную рекомпозицию;
  \item Не забывайте key: items(list, key = \{ it.id \}) при добавлении и удалении трогает только нужные элементы списка — без key удаление одной позиции может перерисовать весь список;
  \item Режим Compose в Layout Inspector: в реальном времени показывает, сколько раз рекомпонировался каждый Composable, — сразу видно, кто делает это слишком часто; оптимизируйте по горячим точкам.
\end{itemize}

Кроме того, включите отчёты компилятора Compose (Metrics) — они сгенерируют список «какие компоненты нестабильны». На пути вглубь данные всегда надёжнее ощущений.

\begin{tcbpractice}{Практика: добавьте учётной книге «три улучшения»}
(1) Добавьте в приложение для учёта расходов напоминание «записывать расходы каждый день в 20:00» (WorkManager + канал уведомлений + разрешение на уведомления в Android 13); (2) переведите создание базы данных и API на внедрение зависимостей через Hilt; (3) попробуйте с помощью ИИ-инструментов ускорить реализацию функции «фильтр по месяцу», а затем проверьте результат по трёхуровневому методу приёмки. Сделав эти три вещи, вы придадите своему проекту вид настоящей «регулярной армии» — профессиональный и солидный.

\end{tcbpractice}

\begin{tcbsummary}{Краткие итоги главы}
\begin{itemize}
  \item Чувствительные разрешения: объявление в Manifest + запрос во время работы + деликатное объяснение после отказа
  \item Три правила уведомлений: отправлять только полезное, давать возможность отключить, не докучать; нужны и канал, и разрешение — одно без другого не работает
  \item Для отложенных задач используйте WorkManager — не изобретайте велосипед
  \item Hilt управляет объектами, модульность — кодом; в небольшом solo-проекте не спешите их внедрять
  \item Четыре сценария помощи ИИ: написать код, найти информацию, написать тесты, исправить баг; методология — в части 7 (главы 37–47)
\end{itemize}

\end{tcbsummary}
